\documentclass{article}
\usepackage{graphicx} % Required for inserting images
\usepackage[T1]{fontenc} % Wichtig für Trennung von Wörtern
\usepackage[utf8]{inputenc} % Erlaubt Umlaute im Code
\usepackage[ngerman]{babel}
\usepackage{csquotes}
\usepackage{bbm}
\usepackage{amsmath}
\usepackage{amssymb}
\usepackage{amsthm}
\usepackage{mathtools}
\usepackage{url}
\usepackage{enumitem}
\usepackage[
  colorlinks=true,
  linkcolor=blue,
  citecolor=blue,
  urlcolor=blue,
]{hyperref}
\usepackage{cleveref}

% Package zum verweisen auf Literatur
\usepackage[style=numeric]{biblatex}
\addbibresource{Ausarbeitung_Literatur.bib}

% 1,5-fachen Zeilenabstand einstellen
\usepackage{setspace}
\onehalfspacing

\newcommand{\N}{\mathbbm{N}} % natürliche Zahlen
\newcommand{\Z}{\mathbbm{Z}} % ganze Zahlen
\newcommand{\R}{\mathbb{R}} % reelle Zahlen
\newcommand{\C}{\mathbb{C}} % komplexe Zahlen
\newcommand{\Q}{\mathbb{Q}} % rationale Zahlen

\newcommand{\nbdash}{\mbox{-}} % geschützter Bindestrich



\newtheorem{satz}{Satz}
\newtheorem{definition}[satz]{Definition}
\newtheorem{lemma}[satz]{Lemma}
\newtheorem{beispiel}[satz]{Beispiel}
\newtheorem{bemerkung}[satz]{Bemerkung}



\title{Darstellungen ganzer Zahlen \\ [0.5em] \Large Eine Ausarbeitung zum Vortrag \\ Proseminar Kombinatorik \\ Sommersemester 2026}
\date{\today}
\author{Carl-Philip Deerberg \\ \url{carl-philip.deerberg@uni-bielefeld.de}}



\begin{document}

\maketitle
\thispagestyle{empty}
\newpage

\pagenumbering{roman}
\tableofcontents
\newpage


\pagenumbering{arabic}

\section{Einleitung}
  Diese Ausarbeitung zum Vortrag befasst sich mit Darstellungsformen von Zahlen, insbesondere ganzer Zahlen und basiert auf \cite{hougardy2018}. 
    
  \vspace{\baselineskip}

  Während für Menschen die Menge der ganzen Zahlen zwar eine nicht überschaubar große Menge an Zahlen darstellt, sind die Elemente der Menge jedoch gut darstellbar.
  Nehmen wir ein beliebiges Element der ganzen Zahlen, so ist dieses in unserer schriftlichen Darstellung eindeutig, für jeden verständlich und leicht reproduzierbar.
  Gleich ob es eine Zahl in vorstellbarer Größenordnung wie ``2'' oder ``8'' ist, oder eine tendenziell schwerer vorstellbare Zahl wie ``27439'' ist: diese Zahlen sind alle übersichtlich in Schritform festzuhalten.
    
  Von dieser für uns gewohnten Darstellung ausgehend können wir uns nun die Frage stellen, welche Alternativen es geben kann, die Elemente der ganzen Zahlen darzustellen.
  Insbesondere betrachten wir die Frage, in welcher Form Computer Zahlen darstellen, speichern und verarbeiten.

  Um diese Frage zu beantworten müssen, wir auf die Ebene der Datenverarbeitung in einem Computer blicken.
  Für diese digitale Verarbeitung beliebiger Daten hat ein Computer nur die Möglichkeit, diese in Form von sogenannten Bits zu speichern.
  Jedes dieser Bits kann nur zwei Zustände annehmen; 0 oder 1. 
  Also schauen wir uns im Folgenden an, wie man ausschließlich mit Nullen und Einsen Zahlen darstellen kann.

  Grundsätzliche Definitionen wie die Menge der ganzen oder natürlichen Zahlen werden in dieser Ausarbeitung nicht ausführlich behandelt.
  Dennoch eine Bemerkung: Für die folgenden Darstellungen ist es irrelevant, ob wir von der Definition $\N = \{1, 2, 3, \dots\}$ ausgehen, oder ob wir die alternative Definition $\N = \{0, 1, 2, 3, \dots\}$ verwenden.
  Die hier diskutierten Sätze funktionieren in beiden Fällen, auch wenn wir von der erstgenannten Definition ausgehen.

\newpage

\section{$b$-adische Darstellung natürlicher Zahlen}
  Bevor wir die Frage für beliebige ganze Zahlen beantworten, betrachten wir zunächst eine übliche Darstellungsform für natürliche Zahlen, die sogenannte $b$-adische Darstellung.
  Mathematisch lässt sich diese durch den folgenden Satz beschreiben: 

  \begin{satz}[$b$-adische Darstellung] \label{satz:b-adische_Darstellung}
    Sei $b\in \N,~ b \geq 2$ und $n \in \N$. 
    Dann existieren $\ell \in \N$ und $z_i \in \{0, \dots, b-1\}, {i = 0, \dots, \ell-1}$, welche die Bedingungen $z_{\ell-1} \neq 0$ und 
    \begin{equation}
      n = \sum_{i=0}^{\ell-1}z_ib^i \label{eq:b-adische_Darstellung_n}
    \end{equation}
    erfüllen.
    $\ell$ und $z_i$ sind dabei eindeutig bestimmt. Dann heißt $z_{\ell-1}\dots z_0$ die \textbf{b\nbdash adische Darstellung} von n und man schreibt auch $n={(z_{\ell-1} \dots z_0)}_b$. \\
    Es gilt: 
    \begin{equation}
      \ell-1 = \lfloor \log_b n \rfloor. \label{eq:b-adische_Darstellung_log}
    \end{equation}
  \end{satz}

  Für den Beweis des Satzes benötigen wir das folgende Lemma:

    \begin{lemma} \label{lemma:Eindeutigkeit_b-adische_Darstellung}
      Sei $k \in \N$ und $\ell \in \N$. 
      Definiere $f^\ell \colon \{0, \dots, k-1\}^\ell \to \{0, \dots, k^\ell-1\}$ durch $f^\ell(w) \coloneq \sum_{i=1}^\ell a_i k^{\ell-i}$ für alle $w=a_1 \dots a_\ell \in \{0, \dots, k-1\}^\ell$.
      Dann ist $f^\ell$ wohldefiniert und bijektiv.
    \end{lemma}
      \begin{proof}
        Der Beweis wird an dieser Stelle nicht ausgeführt, siehe dazu \cite[S.~17f,~Lemma~1.19]{hougardy2018}.
      \end{proof}

    In dieser Ausarbeitung werden die Auf- und Abrundungsfunktionen (auch Gaußklammer genannt) an einigen Stellen verwendet.
    Aus diesem Grund hier die kurze, formale Einführung:
    \begin{bemerkung}[Gaußklammer]\label{bem:Gaussklammer}
      Sei $x \in \R$ und $\lfloor \cdot \rfloor \colon \R \to \Z$ die Abrundungsfunktion bzw. $\lceil \cdot \rceil \colon \R \to \Z$ die Aufrundungsfunktion, so dass
      \begin{align*}
        \lfloor x \rfloor &= \max\{k \in \Z \mid k \leq x\}, \\
        \lceil x \rceil &= \min\{k \in \Z \mid k \geq x\}
      \end{align*}
    \end{bemerkung}
    
    \begin{proof}[Beweis von Satz~\ref{satz:b-adische_Darstellung}] 

      Wir beweisen zunächst die Aussage \eqref{eq:b-adische_Darstellung_log}. 

      Falls gilt, dass $n=\sum_{i=0}^{\ell-1} z_i b^i$ mit $z_i \in \{0, \dots, b-1\}$ für $i = 0, \dots, \ell-1$ und $z_{\ell-1} \neq 0$, also $n$ darstellbar als:
      \begin{equation*}
        n = z_{\ell-1}b^{\ell-1} + z_{\ell-2}b^{\ell-2} + \dots + z_0.
      \end{equation*}
      Dann können wir $n$ nach unten abschätzen:
      \begin{equation*} \tag{$\ast$} \label{eq:BewEindeutigkeit_b-adische_Darstellung_abschaetzung}
        n \geq z_{\ell-1}b^{\ell-1} \geq 1 \cdot b^{\ell-1} = b^{\ell-1}.
      \end{equation*}
      Weiter gilt:
      \[n = \sum_{i=0}^{\ell-1} z_ib^i \underbrace{\leq}_{z_i \leq b-1} \sum_{i=0}^{\ell-1} (b-1)b^i = (b-1) \underbrace{\sum_{i=0}^{\ell-1} b^i}_{\mathclap{\text{geometrische Reihe}}} = (b-1) \frac{b^\ell-1}{b-1} = b^\ell-1,\]
      also ist $n \leq b^\ell-1$ und damit $n < b^\ell$.
      Es gilt dann mit der Ungleichung \eqref{eq:BewEindeutigkeit_b-adische_Darstellung_abschaetzung}: 
      \[b^{\ell-1} \leq n < b^\ell.\] 
      Wir wenden den Logarithmus auf die Terme an: 
      \begin{equation*}
        \log_b(b^{\ell-1}) \leq \log_b(n) < \log_b(b^\ell).
      \end{equation*}
      Da nach Definition des Logarithmus $\log_b(b^x) = x$ gilt, folgt
      \begin{equation*}
        \ell-1 \leq \log_b(n) < \ell,
      \end{equation*}
      also mit Bemerkung \ref{bem:Gaussklammer}
      \begin{equation*}
        \ell-1 = \lfloor \log_b n \rfloor.
      \end{equation*}

      \noindent\textbf{\textit{Beweis der Eindeutigkeit}}

      Die Eindeutigkeit folgt aus der Bijektivität der Abbildung $f^\ell$ aus Lemma~\ref{lemma:Eindeutigkeit_b-adische_Darstellung}.
      Die Länge $\ell$ ist gegeben durch $l\ell= \lfloor \log_b(n) \rfloor +1$ für festes $n$. 
      Da $f^\ell$ injektiv ist, gibt es genau eine Folge $(z_i)_{i=0,\dots,\ell-1}$, welche $n=\sum_{i=0}^{\ell-1} z_ib^i$ erfüllt.
      Die Bedingung $z_{\ell-1} \neq 0$ gilt automatisch mit der Wahl von $\ell$, da $n \geq b^{\ell-1}$.

      \noindent\textbf{\textit{Beweis der Existenz}}

      Wir beweisen die Existenz der $b$-adischen Darstellung durch vollständige Induktion über die Länge der Darstellung $\ell(n) \coloneq 1+ \lfloor \log_b(n) \rfloor$. 
      Eine Erläuterung der vollständigen Induktion findet sich in \cite[S.~26]{hougardy2018} und wird an dieser Stelle vorausgesetzt.

      \textbf{Induktionsanfang:} 
      \begin{equation*}
        \ell(n) = 1 = 1+ \lfloor \log_b(n) \rfloor \Rightarrow \lfloor \log_b(n) \rfloor = 0
      \end{equation*}
      Das gilt für alle $n$ mit $b^0 \leq n < b^1 \Rightarrow n \in \{1, \dots, b-1\}$. 
      Dann ist die Darstellung \[n = \sum_{i=0}^0 z_i b^i = z_0 b^0 = z_0\] trivial.

      \textbf{Induktionsschritt: $\ell-1 \to \ell$} 

      Sei $n\in\N$ mit $\ell(n) = \ell \geq 2$, setze $n' \coloneq \left\lfloor \displaystyle\frac{n}{b} \right\rfloor$.
      Für $\ell' = \ell(n')$ gilt \[\ell' = \left\lfloor \log_b \left\lfloor \frac{n}{b} \right\rfloor \right\rfloor + 1 = \left\lfloor \log_b (\frac{n}{b}) \right\rfloor + 1 = \left\lfloor \log_b n \right\rfloor - 1 + 1 = \ell - 1.\]
      Dann existiert nach Induktionsvoraussetzung für $n'$ eine Darstellung 
      \[n' = \sum_{i=0}^{\ell' -1} z'_i b^i\] 
      mit $z'_i \in \{0,\dots,b-1\} \text{ und } z'_{\ell'-1} \neq 0$. \\
      Setze $z_0 \coloneq n \bmod b \in \{0, \dots, b-1 \}$. Dann gilt:
      \[n = b \cdot \left\lfloor \frac{n}{b} \right\rfloor + (n \bmod b) = b \cdot n' + z_0.\]
      Einsetzen von $n' = \displaystyle\sum_{i=0}^{\ell'-1}z'_i b^i$ nach Induktionsvoraussetzung ergibt
      \begin{align*}
        n = b\cdot \left(\sum_{i=0}^{\ell'-1}z'_i b^i\right) + z_0 = \left(\sum_{i=0}^{\ell'-1}z'_i b^{i+1}\right) + z_0
        \underbrace{=}_{\mathclap{j\coloneq i+1}} \left(\sum_{j=1}^{\ell'} z'_{j-1}b^j\right) + z_0.
      \end{align*}
      Da $\ell' = \ell-1$, setze $z_j \coloneq z'_{j-1}$ für $j \in \{1,\dots,\ell-1\}$ und wir erhalten die gewünschte Darstellung
      \begin{equation*}
        n = \sum_{j=1}^{\ell-1} z_j b^j + z_0 b^0 = \sum_{j=0}^{\ell-1} z_jb^j
      \end{equation*}

    \end{proof}

  \begin{beispiel} \label{bsp:Binärdarstellung}
    Die Zahl $106$ lässt sich, ähnlich zur üblichen Dezimaldarstellung, $b$-adisch wie folgt schreiben:
    \begin{align*}
      106 &= 1 \cdot 10^2 + 0 \cdot 10^1 + 6 \cdot 10^0 \\
          &= 1 \cdot 2^6 + 1 \cdot 2^5 + 0 \cdot 2^4 + 1 \cdot 2^3 + 0 \cdot 2^2 + 1 \cdot 2^1 + 0 \cdot 2^0
    \end{align*}
    bzw. dann in für Computer verarbeitbarer Binärdarstellung $1101010$. 
    Wie wir in Satz~\ref{satz:b-adische_Darstellung} gesehen haben, kann man jede beliebige Zahl $b\geq 2$ als Basis für die Darstellung der natürlichen Zahl nehmen.
  \end{beispiel}

  \vspace{\baselineskip}

  Die Rechenoperationen innerhalb der natürlichen Zahlen funktionieren in der Binärdarstellung ähnlich der üblichen Operationen wie im gewöhnlichen System.

  \vspace{\baselineskip}

  Folgendes sei noch kurz anzumerken:
  Die Wahl von $b$ mit der entsprechenden Einschränkung hat auch eine praktische Relevanz für die Datenverarbeitung.
  Computer nutzen außer der bereits diskutierten Binärdarstellung (zur Basis $b=2$) die Oktaldarstellung ($b=8$) und auch die Hexadezimaldarstellung ($b=16$).

\newpage


\section{$b$-Komplementdarstellung ganzer Zahlen}

  Nachdem wir uns die Darstellung der natürlichen Zahlen angeschaut haben, stellen wir uns die Frage, wie wir die ganzen Zahlen entsprechend darstellen können.
  Bei unserer gewohnten Darstellung besteht der Unterschied zwischen den natürlichen Zahlen und den ganzen Zahlen im wesentlichen lediglich in dem vorangestellten Vorzeichen:
  Für die ganzen Zahlen stellen wir jedem Element der natürlichen Zahlen ein ``$-$'' voran und bezeichnen die Menge als $\Z = \{\dots,-3, -2, -1, 0, 1,2,3,\dots\}$

  Lässt sich dieses System auch analog auf eine für Computer verarbeitbare Darstellung übertragen?
  Schaut man sich das oben betrachtete System aus Bits, welche die Zustände $0$ und $1$ annehmen können, wäre es eine Möglichkeit, das erste Bit für das jeweilige Vorzeichen zu reservieren.
  Man spricht dann von einer sogenannten Vorzeichen- oder auch Betragsdarstellung von Zahlen.


  \begin{beispiel} \label{bsp:Binärdarstellung_Vorzeichen}
    Eine einfache Darstellung ist etwa die Binärdarstellung mit vier Bits.
    Das erste Bit verwenden wir, um das Vorzeichen darzustellen.
    Mit den weiteren drei Bits wird dann der Betrag der Zahl dargestellt.
    Die $0$ im ersten Bit steht dabei für ein positives Vorzeichen, die $1$ für ein negatives.
    So entspricht $+5$ $0101$ und $-5$ entspricht $1101$.
    Durch diese Darstellung lassen sich die Zahlen $-7$ bis $+7$ wie folgt darstellen:
    \begin{align*}
      0 \text{ entspricht } 0000  ~&;~  -0 \text{ entspricht } 1000 \\
      1 \text{ entspricht } 0001  ~&;~  -1 \text{ entspricht } 1001 \\
      & \vdots \\
      7 \text{ entspricht } 0111  ~&;~  -7 \text{ entspricht } 1111
    \end{align*}
  \end{beispiel}

  \vspace{\baselineskip}
  
  Ein Nachteil ist offensichtlich: Es gibt zwei verschiedene Darstellungen für die $0$.
  Während dies in der Datenverarbeitung vernachlässigbar wäre, ergibt sich ein Problem bei der Rechenoperation der Addition.
  Bei der Binärdarstellung für natürliche Zahlen, vergleiche Beispiel~\ref{bsp:Binärdarstellung}, funktioniert die Addition Bitweise wie bei der gewöhnlichen, schriftlichen Addition spaltenweise.
  Dies ist analog zu der gewöhnlichen schriftlichen Addition, die auch spaltenweise abläuft.
  Dort ergibt beispielsweise $0001 + 0001 = 0010$, also $1+1=2$. 
  Dementsprechend muss im Beispiel~\ref{bsp:Binärdarstellung_Vorzeichen} dann $0010+1001=0001$, also $2+(-1)=1$, ergeben.

  \vspace{\baselineskip}

  Aufgrund dieser Schwierigkeit wird im Regelfall eine andere Darstellung für die ganzen Zahlen verwendet.

  \subsection{$b$-Komplementdarstellung - Einführung und Beispiel}

    Man nimmt jede negative Zahl $z\in \{-2^{\ell-1},\dots,-1\}$ mit $z+2^\ell$, wobei $\ell$ hier für die Anzahl der verwendeten Bits steht. 
    Da jedes Bit zwei Zustände annehmen kann, lassen sich $2^\ell$ Zahlen mit $\ell$ Bits darstellen.
    Da auch bei dieser Darstellung prinzipiell ein Bit für das Vorzeichen reserviert wird, ist dieses Bit nicht zur Darstellung des Betrags der Zahl verfügbar.
    Dementsprechend reduziert sich die Anzahl der darstellbaren Zahlen mit $\ell$ Bits auf $2^{\ell-1}$.

    \begin{beispiel} \label{bsp:Komplementdarstellung_Erklärung}
      Bei $\ell=4$ verfügbaren Bits wird die Zahl $z=-3$ wie folgt dargestellt:
      \begin{align*}
        z+2^\ell = -3+16 = 13
      \end{align*}
      $13$ wird in Binärform zu $1101$.
      Da das erste Bit eine $1$ ist, hat dieses einen negativen Wert und man interpretiert die Bits wie folgt:
      \begin{align*}
        &\text{Bit 3: } -2^3 = -8 \\
        &\text{Bit 2: } 2^2 = 4 \\
        &\text{Bit 1: } 2^1 = 2 \\
        &\text{Bit 0: } 2^0 = 1
      \end{align*}
      Konkret für $1101$ bedeutet dies
      \begin{equation*}
        (1 \cdot (-8)) + (1 \cdot 4) + (0 \cdot 2) + (1 \cdot 2) = -3
      \end{equation*}
    \end{beispiel}

    \vspace{\baselineskip}

    Also lassen sich mit vier Bits die folgenden Zahlen darstellen:
    \begin{beispiel} \label{bsp:Komplementdarstellung_4Bits}
      \begin{align*}
        0 \text{ entspricht } 0000  ~&;~  -8 \text{ entspricht } 1000 \\
        1 \text{ entspricht } 0001  ~&;~  -7 \text{ entspricht } 1001 \\
        & \vdots \\
        7 \text{ entspricht } 0111  ~&;~  -1 \text{ entspricht } 1111
      \end{align*}    
      Die Addition funktioniert wie bereits beschrieben:
      \begin{align*}
        &0010 \text{ also } 2 \\
        +~ &1001 \text{ also } -7 \\
        =~ &1011 \text{ also } -5
      \end{align*}
    \end{beispiel}


\subsection{$b$-Komplementdarstellung - allgemeine Formulierung}

  Wir wollen die $b$-Komplementdarstellung einer Zahl nun allgemein und für beliebige Basen $b \geq 2$ beschreiben.

  Zunächst definieren wir hierfür das $b$-Komplement in allgemeiner Form:

  \begin{definition}[$\ell$ stelliges $b$-Komplement] \label{def:b-Komplement}
    Seien $\ell,b \in \N$, mit $b\geq 2$, sei $n \in \{0, \dots, b^\ell-1\}$.
    Dann ist 
    \[K^\ell_b(n) \coloneq (b^\ell-n) \bmod b^\ell\]
    das \textbf{$\ell$-stellige $b$-Komplement} von $n$.
    Für die Basis $b=2$ nennt man $K_2$ das Zweierkomplement, für $b=10$ nennt man $K_{10}$ das Zehnerkomplement.
  \end{definition}
  
  Weiter betrachten wir, wie man im allgemeinen Fall das Komplement einer Zahl berechnet.

  \begin{lemma}[Berechnen des $b$-Komplements] \label{lem:b-Komplement_berechnen}
    Seien $\ell,b \in \N$, mit $b\geq 2$. Sei $n=\displaystyle\sum_{i=0}^{\ell-1} z_ib^i$ mit $z_i \in \{0, \dots, b-1\}$ für $i = 0, \dots, \ell-1$.
    Dann gilt:
    \begin{enumerate}[label=(\roman*)]
      \item $K^\ell_b(n+1) = \displaystyle\sum_{i=0}^{\ell-1}(b-1-z_i)b^i, \text{ falls } n\neq b^\ell-1$ \label{lem:b-Komplement_berechnen_Berechnung}
      \item $K^\ell_b(0)=0$ \label{lem:b-Komplement_berechnen_von0} 
      \item $K^\ell_b(K^\ell_b(n))=n$ \label{lem:b-Komplement_berechnen_KomplementKomplement} 
    \end{enumerate}
  \end{lemma}

    \begin{proof}
      \begin{description}
        \item[\ref{lem:b-Komplement_berechnen_von0}.] Folgt aus Definition~\ref{def:b-Komplement}: $K_\ell^b(n) = (b^\ell-n) \bmod b^\ell$, mit $n=0$:
          \[K_\ell^b(0)=(b^\ell-0) \bmod b^\ell = b^\ell \bmod b^\ell = 0\]
        \item[\ref{lem:b-Komplement_berechnen_Berechnung}] Sei nun $n \in \{0,\dots,b^\ell-2\}$.
          Es gilt nach Definition~\ref{def:b-Komplement} für $m \in \{0, \dots, b^\ell-1\}:$ 
          \[K^b_\ell(m)=(b^\ell-m) \bmod b^\ell.\]
          Da $0 \leq n \leq b^\ell -2$ ist, folgt für $m = n+1$:
          \[1 \leq n+1 \leq b^\ell-1.\]
          Das heißt dann gilt für $(b^\ell - (n+1))$ 
          \[1 \leq b^\ell-(n+1) \leq b^\ell -1,\] 
          liegt also in $\{1, \dots, b^\ell-1\}$.\\
          Es folgt, dass die Modulo-Operation trivial ist:
          \[K^b_\ell(n+1)=(b^\ell-(n+1)) \bmod b^\ell =b^\ell-1-n.\]
          Wegen der geometrischen Reihe $b^\ell-1 = \displaystyle\sum_{i=1}^{\ell-1} (b-1)b^i$ gilt:
          \begin{align*}
            K^b_\ell(n+1) &= b^\ell-1-n \\
            &= \left(\sum_{i=0}^{\ell-1}(b-1)b^i\right) -n \\
          &\underbrace{=}_{\mathclap{\text{Satz \ref{satz:b-adische_Darstellung}}}} \left(\sum_{i=0}^{\ell-1}(b-1)b^i\right) - \underbrace{\left(\sum_{i=0}^{\ell-1} z_ib^i\right)}_{\mathclap{b\text{-adisch}}} \\
            &= \sum_{i=0}^{\ell-1} ((b-1)b^i - z_ib^i) \\
            &= \sum_{i=0}^{\ell-1} (b-1-z_i) b^i
          \end{align*}
        \item[\ref{lem:b-Komplement_berechnen_KomplementKomplement}] Für $n=0$ gilt nach \ref{lem:b-Komplement_berechnen_von0}: 
          \[K^b_\ell(K^b_\ell(0)) = K^b_\ell(0) = 0\]
          Sei nun $n \in \{1, \dots, b^\ell-1\}$ und wir setzen $m \coloneq K_\ell^b(n)$.
          Wegen der Wahl von $n$ gilt dann 
          \[m = (b^\ell-n) \bmod b^\ell = b^\ell-n.\]
          Wir wenden $K_\ell^b$ an:
          \[K_\ell^b(m) = (b^\ell-m) \bmod b^\ell\]
          und substituieren $m=b^\ell-n$:
          \[K_\ell^b(m)=(b^\ell-(b^\ell-n)) \bmod b^\ell = n \bmod b^\ell.\]
          Da $0 < n < b^\ell$, gilt $n \bmod b^\ell = n$ und es folgt
          \[K^b_\ell(K^b_\ell(n))=n.\]
      \end{description}
      Somit ist Lemma \ref{lem:b-Komplement_berechnen} bewiesen.
    \end{proof}

  \vspace{\baselineskip}
  Betrachten wir das folgende Beispiel und nutzen Lemma~\ref{lem:b-Komplement_berechnen}~\ref{lem:b-Komplement_berechnen_Berechnung}, um das $b$-Komplement von positiven Zahlen zu berechnen.
  Prinzipiell subtrahiert man erst $1$ von der ursprünglichen Zahl, um dann stellenweise die Differenz zu $b-1$ zu berechnen:

  \begin{beispiel}
    \begin{enumerate}
      \item Sei $b=10, \ell=4$. Wir berechnen $K_4^{10}(4809)$.
        \begin{enumerate}
          \item $n$ bestimmen: $n+1 = 4809 \Rightarrow n = 4809 -1 = 4808$
          \item $n$ in Ziffern $z_i$ zerlegen: $n=4 \cdot 10^3 + 8 \cdot 10^2 + 0\cdot 10^1 + 8\cdot 10^0$. \\
           Also: $z_3 = 4, z_2 = 8, z_1 = 0, z_0 =8$.
          \item Einsetzen ergibt:
            \begin{align*}
              K_4^{10}(4808+1) &= \sum_{i=0}^3(10-1-z_i)10^i \\
              &= (9-4) \cdot 10^3 + (9-8) \cdot 10^2 \\ 
              &~~+ (9-0) \cdot 10^1 + (9-8) \cdot 10^0 \\
              &= 5 \cdot 10^3 + 1 \cdot 10^2 + 9 \cdot 10^1 + 1 \cdot 10^0.
            \end{align*}
        \end{enumerate}
        Also ist $K_4^{10}(4809)=(5191)_{10}$.
      \item Sei $b=2, \ell=4$. Wir berechnen $K_4^2(6)$. $6$ entspricht $(0110)_2$.
        \begin{enumerate}
          \item $n$ bestimmen: $n+1 = (0110)_2 \Rightarrow n = (0110)_2 -1 = (0101)_2$
          \item $n$ in Ziffern $z_i$ zerlegen: $n=0 \cdot 2^3 + 1 \cdot 2^2 + 0\cdot 2^1 + 1\cdot 2^0$. \\
           Also: $z_3 = 0, z_2 = 1, z_1 = 0, z_0 =1$.
          \item Einsetzen ergibt:
            \begin{align*}
              K_4^2((0110)_2) &= \sum_{i=0}^3(2-1-z_i)2^i = \sum_{i=0}^3(1-z_i)2^i \\
              &= (1-0) \cdot 2^3 + (1-1) \cdot 2^2 + (1-0) \cdot 2^1 + (1-1) \cdot 2^0 \\
              &= 1 \cdot 2^3 + 0 \cdot 2^2 + 1\cdot 2^1 + 0\cdot 2^0 \\
              &= (1010)_2
            \end{align*}
        \end{enumerate}
        Also ist $K_4^2((0110)_2) = (1010)_2$.
    \end{enumerate}
  \end{beispiel}

  \vspace{\baselineskip}

  Nun können wir ganze Zahlen in der $b$-Komplementdarstellung darstellen und speichern.
  Damit Computer adäquat damit rechnen können, ist es vorteilhaft, wenn diese Zahlen dabei eine einheitliche Struktur aufweisen.
  Konkret geht es darum, beliebige positive oder negative ganze Zahlen in einem festgelegten Speicherbereich von $l$ Stellen unterbringen zu können.

\subsection{$\ell$-stellige $b$-Komplementdarstellung}
  Möglich gemacht wird diese einheitliche Darstellung mit der folgenden Definition:

  \begin{definition}[$\ell$-stellige $b$-Komplementdarstellung] \label{def:l-stellige_b-Komplementdarstellung}
    Seien $b,\ell \in \N$ mit $b \geq 2$ und sei $n \in \left\{- \left\lfloor \displaystyle\frac{b^l}{2} \right\rfloor, \dots, \left\lceil \displaystyle\frac{b^\ell}{2} \right\rceil -1 \right\}$. 
    Man erhält die \textbf{$\ell$-stellige $b$-Komple-mentdarstellung} von $n$, indem man die \hyperref[satz:b-adische_Darstellung]{$b$-adische Darstellung} von $n$, falls $n \geq 0$,
    bzw. die $b$-adische Darstellung von $K_b^\ell(-n)$, falls $n < 0$, von der vordersten Stelle mit so vielen Nullen auffüllt, bis sich eine $\ell$-stellige Zahl ergibt.
  \end{definition}

  \vspace{\baselineskip}

  Mit dieser Darstellungsform kann der Computer alle Zahlen einheitlich speichern.
  Darüber hinaus ist der Vorteil der $b$-Komplementdarstellung, dass wir bzw. der Computer mit dieser ohne Fallunterscheidungen rechnen können.
  Es ist also irrelevant, ob wir positive oder negative ganze Zahlen vorliegen haben.
  Einfach gesprochen können wir in der $b$-Komplementdarstellung Addition und Subtraktion so durchführen, als ob wir immer positive ganze Zahlen vorliegen hätten.

  Der folgende Satz begründet dabei, dass sich alle negativen ganzen Zahlen so in positive Zahlen im Sinne der $b$-Komplementdarstellung umwandeln lassen, um wie oben beschrieben rechnen zu können:

  \begin{satz} \label{satz:Rechnen-Komplementdarstellung}
    Seien $b,\ell \in \N$ mit $b \geq 2$. Setze $Z \coloneq \left\{ - \left\lfloor \displaystyle\frac{b^\ell}{2} \right\rfloor , \dots, \left\lceil \displaystyle\frac{b^\ell}{2} \right\rceil -1 \right\}$.\\
    Definiere die Funktion $f \colon Z \to \{0, \dots, b^\ell-1\}, z \mapsto 
      \begin{cases}
        z &\text{falls } z \geq 0 \\
        K_\ell^b(-z) &\text{falls } z <0
      \end{cases}$.\\
    Dann ist $f$ bijektiv und es gilt für $x,y\in Z$:
    \begin{enumerate}[label=(\roman*)] 
      \item Ist $x+y \in Z$, so gilt 
        \[f(x+y) = (f(x) + f(y)) \bmod b^\ell.\] \label{satz:Rechnen-Komplementdarstellung_Addition}
      \item Ist $x \cdot y \in Z$, so gilt 
        \[f(x \cdot y) = (f(x) \cdot f(y)) \bmod b^\ell.\] \label{satz:Rechnen-Komplementdarstellung_Multiplikation}
    \end{enumerate}
  \end{satz}

    \begin{proof}
      \begin{description}
        \item[Bijektivität:] Wir betrachten Definitions- und Zielbereich von $f$.
          Die Menge \[Z = \left\{ - \left\lfloor \displaystyle\frac{b^\ell}{2} \right\rfloor , \dots, \left\lceil \displaystyle\frac{b^\ell}{2} \right\rceil -1 \right\}\] enthält genau $b^\ell$ Elemente.
          Die Zielmenge \[\{0,\dots,b^\ell-1\}\] enthält ebenfalls genau $b^\ell$ Elemente.
          Wir zeigen die Injektivität von $f$: 

          Seien $z_1, z_2 \in Z$. Angenommen, $f(z_1) = f(z_2)$.
          Da $f(z) \equiv z (\bmod b^\ell)$, gilt $z_1 \equiv z_2 (\bmod b^\ell)$.
          Da die Differenz von zwei verschiedenen Elementen aus $Z$ im Betrag kleiner als $b^\ell$ sein muss, kann die Bedingung nur gelten, wenn $z_1 = z_2$.
          Es folgt: $f$ ist injektiv.

          \noindent Da $f$ injektiv ist, und $|Z| = b^\ell = |\{0,\dots,b^\ell-1\}|$, also Definitions- und Zielbereich die gleiche Kardinalität haben, ist $f$ bijektiv.
      \end{description}

      \vspace{\baselineskip}

      Seien $x,y \in Z$ und $p,q \in \{0,1\}$. 
      Wir betrachten $f(x)=x + p\cdot b^\ell$ und $f(y)=y + q \cdot b^\ell$.
      
      \begin{description}
        \item[\ref{satz:Rechnen-Komplementdarstellung_Addition}] Behauptung: $f(x+y) = (f(x)+f(y)) \bmod b^\ell$ gilt, wenn $(x+y) \in Z$. 
          
          Es gilt:
          \begin{align*}
            (f(x)+f(y)) \bmod {b^\ell} &= ((x+p\cdot b^\ell) + (y+q\cdot b^\ell)) \bmod {b^\ell} \\
            &= (x+y+p \cdot b^\ell + q \cdot b^\ell) \bmod {b^\ell} \\
            &= (x+y+ \underbrace{(p+q)\cdot b^\ell}_{\mathclap{\text{Vielfaches von }b^\ell}}) \bmod {b^\ell} \\
            &= (x+y) \bmod {b^\ell} \\
            &= (x+y+b^\ell) \bmod {b^\ell} \\
            &= f(x+y).
          \end{align*}

        \item[\ref{satz:Rechnen-Komplementdarstellung_Multiplikation}] Behauptung: $f(x \cdot y) = (f(x) \cdot f(y)) \bmod b^\ell$ gilt, wenn $(x\cdot y) \in Z$.
          
          Es gilt:
          \begin{align*}
            f(x+y) &= (x+y+b^\ell) \bmod b^\ell \\
            &= (x+y) \bmod b^\ell \\
            &= (x+y+0) \bmod b^\ell \\
            &= (x+y+ \underbrace{(p+1)\cdot b^\ell}_{\mathclap{\equiv 0 (\bmod b^\ell)}}) \bmod b^\ell \\
            &= (x+ p\cdot b^\ell + y + q\cdot b^\ell) \bmod b^\ell \\
            &= ((x + p \cdot b^\ell) + (y + q\cdot b^\ell)) \bmod b^\ell \\
            &= (f(x) + f(y)) \bmod b^\ell.
          \end{align*}
      \end{description}
    \end{proof}

  \vspace{\baselineskip}

  Mit den hier eingeführten Darstellungen für natürliche und ganze Zahlen haben wir Möglichkeiten gesehen, die Zahlen auf alternative Art und Weise darzustellen.
  Wie bereits an mehreren Stellen erwähnt, bieten diese Formen besonders für die Verarbeitung in Computern Vorteile.
  Aufgrund der diskutierten Sätze und Eigenschaften können wir die gewöhnlichen Rechenoperationen nun auch innerhalb der $b$-Komplementdarstellung von ganzen Zahlen anwenden, ohne auf Ungenauigkeiten oder Probleme bei der Verarbeitung zu stoßen.

\newpage
\printbibliography[heading=bibintoc]

\end{document}